\documentclass[12pt]{article}

\usepackage{graphicx}
\usepackage{url}
\usepackage{hyperref}
\usepackage{verbatim}
\usepackage[title,titletoc,toc]{appendix}
\usepackage{wasysym}

\renewcommand{\topfraction}{1.0}
\renewcommand{\bottomfraction}{1.0}
\renewcommand{\textfraction}{0.0}
\renewcommand{\floatpagefraction}{1.0}
\renewcommand{\dbltopfraction}{1.0}


\begin{document}

\begin{titlepage}

\begin{center}
{\large STAT 5810/6810 Introduction to Statistical Computing} \\[4cm]

{\LARGE \bf Solutions for Homework  3} \\[1cm]
by \\[0.5cm]
{\bf J\"urgen Symanzik} \\[2.5cm]
{\bf Date:} \today \\[2cm]

UTAH STATE UNIVERSITY \\[0.5cm]
Logan, UT \\[0.5cm]
Fall 2012 \\[0.5cm]
\end{center}

\thispagestyle{empty}
\vfill
\end{titlepage}


\newpage 


\pagenumbering{roman}

\tableofcontents


\newpage


%\listoftables
%\addcontentsline{toc}{section}{List of Tables}
%
%\newpage

\listoffigures
\addcontentsline{toc}{section}{List of Figures}

\newpage 


\pagenumbering{arabic}


\section{Solutions for Homework 3}

\subsection{General Comments}

Here are my solutions for Homework 3 --- and, as usual, a few more \LaTeX\ features! 
In particular, note how we can use the {\tt appendix} \LaTeX\ package, rather than
the \verb|\appendix| command. Noteworthy are the numbering and the modified
entries in the table of contents.

Specific results and figures for HW03 can be found in Section~\ref{Results}.
R code for all questions can be found in Appendix~\ref{AppendixWithRCode}.


\subsection{Self--Explanatory Figures}

When creating figures and figure captions, it is expected that these are
self--explanatory. This means that a reader gets sufficient information
by just looking at a figure and its caption without having to read the
main text. The main text is used to provide additional details beyond
the information that is stated in the figure caption. Here are some
useful links:
\begin{itemize}
\item \url{http://unilearning.uow.edu.au/report/1fi.html}
\item \url{http://ocw.mit.edu/courses/biological-engineering/20-109-laboratory-fundamentals-in-biological-engineering-fall-2007/assignments/guide_lab_report-htm/}
\end{itemize}

By the way, what holds for figures usually also holds for tables.
google yourself for addional details.


\newpage


\section{Results}\label{Results}

In this section, we provide general results and include the four figures from HW03. 


\subsection{World Record in Men's 1500 Meter Run}

In our first plot (Figure~\ref{HW03_Plot1}), we examine the development
of the world record times in men's 1500 meter run from 1892 to 2012.
Overall, the time for the world record  was reduced from about 265 sec to about 206 sec,
i.e., roughly by a minute, during the last 120 years.
The world record was broken (or equalized) 51 times by 41 different
world record holders.

\begin{figure}[h]
\centerline{\includegraphics[width=0.99\textwidth]{HW03_Plot1.pdf}}
\caption{\label{HW03_Plot1}
This figure shows the world record times (in sec) in men's 1500 meter run 
from 1892 to 2012. Marked are two periods (1944 and 1998) when the
world record wasn't further improved for ten and fourteen years,
respectively. The names of the two world record holders during these
periods are indicated. 
}
\end{figure}


\newpage


\subsection{Total Medals in Relation to Population and GDP}

Figure~\ref{HW03_Plot2} shows a three--way relationship between the
population of a country (horizontal), the GDP per person (vertical),
and the 2012 total medals won (area of the circles) for those
countries that won at least one medal at the 2012 Summer Olympics
in London. While there is no noticeable relationship between
GDP per person and population (both on the log scale),
it is obvious that countries with the largest total number 
of medals have a relatively large population and a 
relatively high GDP per person.

As a side note, recall that the HW instructions asked you to
{\it ``create circles for the plotting
symbols where the area of the circle is proportional to the 
total number of medals.''}
The R help page for {\tt symbols} specifies:
\begin{verbatim}
      circles        a vector giving the radii of the circles.
\end{verbatim}
Thus, we have to use the {\bf square root} of the total number of medals
in the R code, and not just the total number of medals itself.
As a reminder, the area $A$ of a circle is calculated as $A = \pi \cdot r^2$,
where $r$ is the radius.\footnote{I suppose you know this formula,
but I wanted to use this opportunity to introduce
mathematical formulas and the {\tt footnote} command \smiley.
And yes, \smiley\ and \frownie\
can be obtained from the \LaTeX\ package {\tt wasysym}.
}


\newpage


\begin{figure}[h]
\centerline{\includegraphics[width=5in]{HW03_Plot2.pdf}}
\caption{\label{HW03_Plot2}
This figure shows a scatterplot of the
2012 total medals in relation to population and GDP per person.
Both axes are drawn on a log--scale.
The area of the circles is proportional to the total number of medals won
for each country. Highlighted (in blue) are the country codes for
those five countries with the largest number of medals won.
Also highlighted (in red) is my home country, Germany (``DEU''). 
}
\end{figure}


\newpage


\subsection{Map of Total Medals by Country}

Figure~\ref{HW03_Plot3} shows a map of the world where the number of medals won
by country has been overlaid as circles. Some inconsistencies become immediately
apparent (when the circles are small enough). First of all, there are numerous circles
in the area of the former USSR. The same is true when we zoom into south--east
Europe where we notice that there are also numerous circles in the area of the former Yugoslavia.

A comparison of names used in the map with names used in the data file is helpful:
\begin{verbatim}
   mapCountries = map(namesonly = TRUE)
   SO2012Ctry$Country[!(SO2012Ctry$Country %in% mapCountries)]
\end{verbatim}

Overall, there are 85 country names that do not match. This can be expected for
names with slightly different spellings, in particular when spaces are removed
from the country names, such as UnitedArabEmirates, UnitedKingdom, and NorthKorea.
Moreover, we see that the countries of the former USSR (from Armenia to Uzbekistan)
and of the former Yugoslavia (from Bosnia to Slovenia) appear in the data file,
but not in the map.

This outcome should not have come totally unexpected. The documentation of the
{\tt maps} package at \url{http://cran.r-project.org/web/packages/maps/} states:
{\it ``Original S code by Richard A. Becker and Allan R. Wilks.''} So, this
is based on a map from the 1980s. A more recent version of a world map is
accessible via the R package {\tt maptools} (see
\url{http://cran.r-project.org/web/packages/maptools/}
for details).

\underline{\bf Note:} As a side note, I rescale the width of Figure~\ref{HW03_Plot3}
to \verb|1.5\textwidth|. Typically, we only rescale to a maximum
of 0.99 times the textwidth so that a figure occupies the maximum
possible width on a page. Here, the original pdf figure produced in R
contains a lot of white space on both sides of the figure. So, after rescaling,
the figure actually extends into the margins on both sides, but this
is fortunately not visible.


\newpage


\begin{figure}[t]
\centerline{\includegraphics[width=1.5\textwidth]{HW03_Plot3.pdf}}
\caption{\label{HW03_Plot3}
This figure shows the map of the world, with the 2012 total number of
medals by country overlaid as circles. The area of the circles is proportional 
to the total number of medals for each country. Some inconsistencies
exist and are further discussed in the text.
}
\end{figure} ~\\


\newpage


\subsection{Participants by Sex and Sports}

\begin{figure}[b]
\centerline{\includegraphics[width=4.5in]{HW03_Plot4.pdf}}
\caption{\label{HW03_Plot4}
Side--by--side bar charts of female and male participants in the 2012
Olympics by sport. The disciplines are sorted increasingly with respect to the
total number of participants of both genders in each sport. 
}
\end{figure}

Figure~\ref{HW03_Plot4} shows side--by--side bar charts of the female and
male participants in the 2012 Olympics. By far the most participants (of both genders)
participate in athletics. The next four most popular disciplines are
swimming, football ($=$ soccer), rowing, and cycling. The numbers of
female and male participants in most sports are similar, possibly with
a few more men participating. However, four sports have a rather different
gender distribution: wrestling and boxing have mostly male participants,
while synchronised swimming and gymnastics rhythmics have only
female participants.

\underline{\bf Note:} There is a problem in the data. Did you notice
on the lower left side of Figure~\ref{HW03_Plot4}  that there is
a single participant in cycling --- but that cycling is also the fifth
most popular sport? A closer look at the actual data resolves this problem:
One participant's sport is listed as {\it ``Cycling~''} (with an extra space 
at the end), rather than
{\it ``Cycling''}. This should be corrected, either in the data file
or in R, and a corrected plot should be produced.


\newpage


\begin{appendices}


\section{R Code for Homework 3}\label{AppendixWithRCode}

Here is the R code for Homework 3:

\scriptsize
\verbatiminput{HW03_Graphics_sol.R}

\end{appendices}


\end{document}

